\section{数据与计算：ViT 何时超过卷积基线}

架构图本身无法回答 ViT 是否有效。讲座接着比较不同 pretraining dataset 与 compute budget，说明 ViT 在小数据上可能不如强正则化 CNN，却能在更大数据和合适 recipe 下取得更好的 transfer。这里的因果对象是完整训练系统，而不是 attention 单组件。

\subsection{Scaling with data：弱先验需要更多证据}

当预训练从 ImageNet 扩展到 ImageNet-21k 与 JFT-300M，ViT 的性能提升更明显。直觉上，强卷积先验在低数据区间提供 sample efficiency；数据规模扩大后，较灵活的模型可以学习更适合任务分布的结构，先验不足的代价逐渐被数据抵消。

\lecturefigure{slide-19-scaling-with-data.jpg}{ViT 与 BiT 在不同预训练数据规模下的 ImageNet 与 linear few-shot 结果。}{官方视频 30:24；Dosovitskiy et al. 2020。}

\begin{knowledgebox}{读图：先看横轴，再看交叉点}
左图横轴是预训练数据集规模，纵轴是 ImageNet top-1；右图横轴是 JFT 样本数量，纵轴是 linear few-shot accuracy。小数据端 BiT/ResNet 更稳，大数据端 ViT 曲线上升更快。关键不是某个点领先，而是模型 family 对数据规模的响应斜率不同。曲线不能证明无限数据下永远保持同一趋势。
\end{knowledgebox}

\begin{warningbox}{数据规模不是“图片数量”一个数字}
有效规模还取决于 label noise、class diversity、long-tail coverage、duplicate rate、resolution、domain breadth 与 augmentation。300M 个高重复网页样本不等于 300M 个独立信息单元。比较 scaling curve 时应报告清洗与采样策略。
\end{warningbox}

\subsection{Pretraining frontier：论文结果是系统共同产物}

模型在历史 performance frontier 上的位置受 dataset、hardware、optimizer、training duration 与 hyperparameter search 共同影响。ViT 点位于 ImageNet 模型演进图的上方，说明该 recipe 有竞争力；但这不是“architecture-only”消融，也不意味着每个团队都能在较小预算复现相同增益。

\lecturefigure{slide-20-vit-pretraining.jpg}{ViT 在 ImageNet accuracy 与预训练规模的历史图上形成新的 frontier 点。}{官方视频 30:32。}

\begin{knowledgebox}{读图：frontier 只比较被画出的方案}
散点图上“更上方”代表更高 accuracy，“更右侧”通常代表更多数据或规模。Frontier 表示在已报告方案中较优，不代表 Pareto 轴完整；若没有标出 energy、latency、memory、annotation cost 与 tuning attempts，图无法回答部署效率或研究成本。
\end{knowledgebox}

\subsection{Position embedding：空间结构是学出来还是写进去}

没有 position embedding 的 self-attention 对 token 排列近似 permutation-equivariant，无法区分同一 patch 集合的空间顺序。ViT 学习一组绝对位置向量；可视化它们的相似度，可以检查模型是否从二维图像中恢复出邻近关系和网格结构。

\lecturefigure{slide-21-position-embeddings.jpg}{学习到的位置编码相似度呈现二维邻域结构，并与不同初始化/共享方式比较。}{官方视频 33:56。}

\begin{knowledgebox}{读图：热图不是“模型理解空间”的证明}
每个位置与其他位置的 embedding similarity 若呈局部平滑，说明训练恢复了二维邻近结构。右侧表格比较不同 position embedding 设计的性能。该证据只说明表示中存在空间规律，不证明模型获得物体级几何、因果关系或完整 scene understanding；还需 attention pattern、probe 与 downstream behavior 支持。
\end{knowledgebox}

\teachervoice{Q\&A 中 Lucas 回到这张 slide，解释 position embedding 的主要作用是区分位置，而局部邻接结构可在训练中出现。课堂提示：可视化是机制线索，不是因果证明；应结合移除、打乱、插值和分辨率迁移实验。}

\subsection{Compute-normalized comparison}

若只按参数量比较，ViT、Mixer 与 CNN 的实际计算可能差异很大。讲座用 compute 轴观察 transfer performance：在相近 FLOPs 下，模型结构与训练 recipe 决定能否把计算转化为可迁移表示。读图前要把理论运算量与真实系统成本分开：同一 FLOPs 可能对应完全不同的矩阵形状、memory traffic、并行度与通信，因此这里只比较算法层面的 compute efficiency。

\lecturefigure{slide-22-scaling-with-compute.jpg}{不同模型在预训练 compute 与 transfer accuracy 之间形成不同效率前沿。}{官方视频 39:24。}

\begin{knowledgebox}{读图：FLOPs 不是时间，也不是成本}
横轴 approximate compute 便于算法比较，但相同 FLOPs 在不同硬件上可能有不同 latency、memory traffic 与 utilization。Attention、convolution 和 MLP 的 kernel efficiency 不同；分布式通信也未必计入。图能回答“理论计算预算下谁更准确”，不能直接回答“哪个系统更便宜或更快”。
\end{knowledgebox}

\subsection{本章小结}

ViT 的优势出现在数据、regularization 与 compute 足够的 regime。数据 scaling 曲线展示响应斜率，frontier 图展示历史竞争力，position embedding 提供空间结构线索，compute 图则提醒参数量不是唯一成本尺度。

